\begin{center}
Problema J

Zé do Grafo

{\tiny Por André da Cunha Ribeiro}

Timelimit: $1$	
\end{center}			

{\footnotesize
Zé do Grafo, a lenda que ecoa como o Mestre Jedi exterminador de monstros, emerge das sombras como uma faísca ardente de esperança na enigmática cidade entrelaçada nos recônditos do IFGoiano - Campus Rio Verde. Essa cidade, envolta em mistério, desdobra-se perante olhares curiosos como um enigma arquitetônico complexo, com conexões intrincadas que aguçam a mente.

Nesse ambiente intrigante, Zé do Grafo se ergue como um guardião destemido, dotado de poderes transcendentes. Em um dia sereno, enquanto explorava a cidade, ele desvenda um bairro oculto, um santuário de portais secretos que se estendem ao multiverso. Cada portal é adornado com inscrições enigmáticas, revelando as delicadas ligações entre os diversos universos. Uma característica única surge: cada universo conectado é tingido por uma cor distinta, representando a natureza de sua atmosfera. Nota-se que universos vizinhos exibem atmosferas diversas, delineando uma dança cósmica de elementos distintos.

Entretanto, dilemas pairam sobre Zé do Grafo: ele só pode sobreviver em dois tipos específicos de atmosfera, uma limitação impregnada pelas superstições que permeiam a linhagem de mestres Jedi. A narrativa então faz um compasso, direcionando o holofote para um novo protagonista: você. Sua ajuda, inestimável e ancorada na compreensão profunda dos algoritmos e grafos, é convocado para solucionar este enigma vital. A pergunta que se insinua é crucial: há algum portal capaz de permitir que o lendário Zé do Grafo explore os multiversos, sem que isso desafie suas crenças enraizadas na superstição?

\vspace{0.3cm}
\textbf{Entrada}
Cada entrada contém os dados de cada portal. A primeira linha contém dois inteiros $N$, $M$, indicando respectivamente o número de multiversos $(1 \leq N \leq 50)$ e as ligações dos multiversos $(1 \leq M \leq 2000)$.

Cada uma das $M$ linhas seguintes composta de dois inteiros que descreve as ligações dos multiverso.

\vspace{0.3cm}
\textbf{Saída}
O sistema deve retornar ``Sim'' caso Zé do Grafos consiga explorar todo o multiverso. Caso contrário, a resposta será ``Não''.
\begin{table}[h]
    \centering
    \begin{tabular}{|p{7cm}|p{7cm}|}
        \hline
        \begin{center}
            \vspace{-0.3cm}
            \textbf{Exemplos de Entrada}
            \vspace{-0.3cm}
        \end{center} 
         &  
        \begin{center}
            \vspace{-0.3cm}
            \textbf{Exemplos de Saída}
            \vspace{-0.3cm}
        \end{center} \\ \hline
         \texttt{6 7} & \texttt{Sim} \\  
         \texttt{1 2} & \texttt{} \\
         \texttt{1 6} & \texttt{} \\
         \texttt{2 3} & \texttt{} \\
         \texttt{3 4} & \texttt{} \\
         \texttt{3 6} & \texttt{} \\
         \texttt{4 5} & \texttt{} \\         
         \texttt{5 6} & \texttt{} \\
         \hline
         \texttt{6 7} & \texttt{Não} \\  
         \texttt{1 2} & \texttt{} \\
         \texttt{1 3} & \texttt{} \\
         \texttt{2 3} & \texttt{} \\
         \texttt{3 4} & \texttt{} \\
         \texttt{4 5} & \texttt{} \\
         \texttt{4 6} & \texttt{} \\
         \texttt{5 6} & \texttt{} \\
         \hline
         \texttt{3 3} & \texttt{Não} \\  
         \texttt{1 2} & \texttt{} \\
         \texttt{1 3} & \texttt{} \\
         \texttt{2 3} & \texttt{} \\
         \hline
    \end{tabular}
    \caption{Exemplos de entradas e saídas}
    \label{tab:inputJJ}
\end{table}}

\newpage

\begin{center}
\noindent
\lstinputlisting{abobora_23/J/J3.cpp}
\end{center}

